\pnum{5.17}{Electronic configurations and atomic terms}
The Pauli principle allows two electrons of opposite
spin in each spatial orbital. Filling $1s$, $2s$,
and $2p$ subshells in the usual order gives:
\begin{center}\begin{tabular}{lll}\toprule
Element&Configuration&Ground term\\\midrule
H&$1s$&$\,{}^{2}S_{1/2}$\\
He&$1s^2$&$\,{}^{1}S_0$\\
Li&$1s^2\,2s$&$\,{}^{2}S_{1/2}$\\
Be&$1s^2\,2s^2$&$\,{}^{1}S_0$\\
B&$1s^2\,2s^2\,2p$&$\,{}^{2}P_{1/2}$\\
C&$1s^2\,2s^2\,2p^2$&$\,{}^{3}P_0$\\
N&$1s^2\,2s^2\,2p^3$&$\,{}^{4}S_{3/2}$\\
O&$1s^2\,2s^2\,2p^4$&$\,{}^{3}P_2$\\
F&$1s^2\,2s^2\,2p^5$&$\,{}^{2}P_{3/2}$\\
Ne&$1s^2\,2s^2\,2p^6$&$\,{}^{1}S_0$\\\bottomrule
\end{tabular}\end{center}
A term $\,{}^{2S+1}L_J$ specifies
spin $S$, orbital angular momentum $L$ and
total $J$. For H, He, Li, Be, the term values
follow immediately from their open or
closed $s$ shells. Boron's $2p$ electron
allows $\,{}^{2}P_{1/2}$ and $\,{}^{2}P_{3/2}$.
For carbon $2p^2$, antisymmetry allows
$\boxed{{}^{1}S_0,\ {}^{1}D_2,\ {}^{3}P_{0,1,2}}$.
For nitrogen $2p^3$, its allowed LS terms are
$\boxed{{}^{4}S_{3/2},\ {}^{2}D_{3/2,5/2},\
{}^{2}P_{1/2,3/2}}$.
Hund's energy-ordering rules select the
ground states.

\pnum{5.18}{Hund's three ordering rules}
(a) For excited $1s\,n\ell$ helium,
Hund's first rule favors a spin triplet
$S=1$ over a singlet $S=0$.
The required antisymmetric spatial
factor tends to reduce electron repulsion.
(b) Two $2p$ electrons in $L=2,M_L=2$
must both carry $m_\ell=1$; the orbital
top-of-ladder state is spatially symmetric.
It cannot accompany the symmetric spin
triplet, so at maximum $S=1$ carbon
has $L=1$, not $L=2$.
(c) Boron's $2p$ subshell is less than
half filled with $L=1,S=1/2$, hence
Hund's third rule selects
$\boxed{J=|L-S|=1/2}$,
the term $\,{}^{2}P_{1/2}$.
(d) For carbon, maximize $S=1$, use
allowed $L=1$, and choose $J=0$:
$\boxed{{}^{3}P_0}$.
For nitrogen $2p^3$, all three
$m_\ell=-1,0,+1$ orbitals are singly
occupied with parallel spin:
$S=3/2$, $M_L=0$, $L=0$, and
$J=3/2$, giving
$\boxed{{}^{4}S_{3/2}}$.
Hund rules are ordering principles
for LS-coupled many-electron atoms,
not exact solutions of the Coulomb problem.

\pnum{5.19}{Dysprosium's ground term}
A term $\,{}^{2S+1}L_J$ with label
$\,{}^{5}I_8$ has
$2S+1=5$, giving $\boxed{S=2}$.
The letter $I$ represents
$\boxed{L=6}$, and the subscript
gives $\boxed{J=8}$.
A likely dominant configuration is
$\boxed{[\mathrm{Xe}]\,4f^{10}\,6s^2}$.
The closed $6s^2$ pair contributes
zero net $L$ and $S$.
The open $4f^{10}$ shell supports
$S=2,L=6$, and being more than half
filled favors $J=L+S=8$.

\pnum{5.20}{Mean energy in a filled three-dimensional Fermi sphere}
At zero temperature, spin-$1/2$
electrons fill $k\le k_F$. Counting
two spin orientations,
\begin{equation}
N_e=2\frac{V}{(2\pi)^3}4\pi
\int_0^{k_F}k^2\dd k
=\frac{Vk_F^3}{3\pi^2}.
\end{equation}
Their total energy is
\begin{equation}
E_{\rm tot}=2\frac{V}{(2\pi)^3}4\pi
\int_0^{k_F}
\frac{\sh^2k^2}{2m}k^2\dd k
=\frac{V\sh^2k_F^5}{10\pi^2m}.
\end{equation}
Dividing gives
\begin{equation}
\boxed{\frac{E_{\rm tot}}{N_e}
=\frac35\frac{\sh^2k_F^2}{2m}
=\frac35E_F}.
\end{equation}

\pnum{5.21}{Copper: Fermi energy, speed, temperature and pressure}
The copper mass density is
$\rho_m=8.96\times10^3\,{\rm kg/m^3}$,
molar mass $0.0635\,{\rm kg/mol}$,
with one free electron per atom.
The conduction-electron density is
\begin{equation}
n=\frac{\rho_m N_A}{0.0635\,{\rm kg/mol}}
\simeq8.50\times10^{28}\ {\rm m}^{-3}.
\end{equation}
Using
$E_F=\sh^2(3\pi^2n)^{2/3}/(2m_e)$,
\begin{equation}
\boxed{E_F\simeq7.03\ {\rm eV}}.
\end{equation}
The Fermi speed is
$\boxed{v_F=\sqrt{2E_F/m_e}\simeq
1.57\times10^6\,{\rm m/s}\simeq0.0052c}$,
so the nonrelativistic assumption is safe.
The Fermi temperature and degeneracy pressure are
\begin{equation}
\boxed{T_F=E_F/k_B\simeq8.15\times10^4\ {\rm K}},
\quad
\boxed{P=\frac25nE_F\simeq3.83\times10^{10}\ {\rm Pa}}.
\end{equation}
$E_F$ must be converted to Joules
in the pressure calculation.
These are noninteracting zero-temperature
electron-gas estimates.

\pnum{5.22}{Helium-3 Fermi temperature}
A helium-3 atom has spin one-half
and mass $m_3\simeq5.01\times10^{-27}$ kg.
For density $82\,{\rm kg/m^3}$,
the particle density is
$n=82/m_3\simeq1.64\times10^{28}\,{\rm m^{-3}}$.
For spin degeneracy two,
\begin{equation}
\boxed{T_F=
\frac{\sh^2}{2m_3k_B}(3\pi^2n)^{2/3}
\simeq3.1\ {\rm K}}.
\end{equation}
This ideal-gas estimate omits
liquid helium-3 interactions.

\pnum{5.23}{Bulk modulus of the degenerate electron gas}
At fixed electron count $N_e$,
$n=N_e/V$ and
$P=Cn^{5/3}=CN_e^{5/3}V^{-5/3}$.
Therefore
\begin{equation}
\frac{\dd P}{\dd V}=-\frac53\frac{P}{V},
\qquad
\boxed{B=-V\frac{\dd P}{\dd V}
=\frac53P}.
\end{equation}
Using the copper pressure in 5.21,
\begin{equation}
\boxed{B_{\rm Cu}^{\rm ideal}
\simeq6.38\times10^{10}\ {\rm Pa}}.
\end{equation}
The textbook's comparison value is
$1.34\times10^{11}\,{\rm Pa}$;
the discrepancy is expected from
interactions and band structure
neglected by the gas model.
